\documentclass[10pt,a4paper]{ctexart}

\usepackage[top=2cm,bottom=2cm,left=2cm,right=2cm]{geometry}
\usepackage{amsmath,amssymb}
\usepackage{enumitem}
\usepackage{tasks}
\usepackage{tikz}
\usepackage{circuitikz}
\usepackage{adjustbox}

\settasks{label-width=1.2em}

\newcommand{\fillblank}{\underline{\hspace{1.5cm}}}
\usetikzlibrary{positioning,arrows.meta,shapes.geometric,calc}

\setlength{\parskip}{6pt}
\setlength{\parindent}{0pt}

\ctexset{
  section/format = {\Large\bfseries}
}

\begin{document}

\begin{center}
  {\LARGE\bfseries 信号与系统\ 期末试卷（回忆版）}

  \vspace{0.4cm}
  {\large 说明：本卷根据考试回忆整理，具体数值可能与原卷有出入，但考点和题型保持一致。}

  \vspace{0.25cm}
  {\normalsize
  来源：与 24 级同学对答案整理\\[2pt]
  贡献：McFlurry，ShiningQStar，5555}
\end{center}

\vspace{0.3cm}

% ==================== 一、选择题 ====================
\section*{一、选择题}

\begin{enumerate}[label=\arabic*.,leftmargin=*]

\item 已知系统的输入 $e(t)$ 与输出 $r(t)$ 满足关系
  \[
  r(t)=\int_{t-1}^{t+1} e(\tau)\,\mathrm{d}\tau,
  \]
  则该系统具有\fillblank。

  \begin{tasks}(2)
    \task 线性、时不变、非因果
    \task 线性、时不变、因果
    \task 线性、时变、非因果
    \task 非线性、时不变、非因果
  \end{tasks}

\item 已知某信号的频谱为
  \[
  X(\mathrm{j}\omega)=
  \begin{cases}
  1, & 3 < |\omega| < 7,\\[2pt]
  0, & \text{其他},
  \end{cases}
  \]
  则其傅里叶逆变换 $x(t)=$\fillblank。

  \begin{tasks}(2)
    \task $\dfrac{\sin 7t-\sin 3t}{\pi t}$
    \task $\dfrac{\sin 7t+\sin 3t}{\pi t}$
    \task $\dfrac{2}{\pi t}\sin 5t\cos 2t$
    \task $\dfrac{\sin 7t-\sin 3t}{2\pi t}$
  \end{tasks}

\item 已知周期信号 $f(t)$ 的周期 $T=4$，在一个周期 $[-2,\,2]$ 内的波形为：$t\in[-1,0]$ 时由 $0$ 线性上升至 $2$，$t\in[0,1]$ 时由 $2$ 线性下降至 $0$，其余区间为 $0$。则其指数形式傅里叶级数的系数为\fillblank。

  \begin{tasks}(2)
    \task $F_0=\dfrac12,\quad F_n=\dfrac12\,\mathrm{Sa}^2\!\left(\dfrac{n\pi}{4}\right)$
    \task $F_0=1,\quad\; F_n=\dfrac12\,\mathrm{Sa}^2\!\left(\dfrac{n\pi}{2}\right)$
    \task $F_0=\dfrac12,\quad F_n=\dfrac14\,\mathrm{Sa}^2\!\left(\dfrac{n\pi}{4}\right)$
    \task $F_0=1,\quad\; F_n=\mathrm{Sa}^2\!\left(\dfrac{n\pi}{4}\right)$
  \end{tasks}

\item 已知象函数
  \[
  F(s)=\frac{s^2+2s+3}{(s+1)(s^2+4)},
  \]
  则 $f(0_+)$ 和 $f(\infty)$ 分别为\fillblank。

  \begin{tasks}(2)
    \task $f(0_+)=0,\; f(\infty)=0$
    \task $f(0_+)=0,\; f(\infty)=\dfrac34$
    \task $f(0_+)=1,\; f(\infty)=0$
    \task $f(0_+)=1,\; f(\infty)\text{不存在}$
  \end{tasks}

\end{enumerate}

% ==================== 二、填空题 ====================
\section*{二、填空题}

\begin{enumerate}[label=\arabic*.,leftmargin=*]

\item 已知
  \[
  F(s)=\frac{2s+3}{(s+2)(s+1)},\qquad \text{收敛域：}-2<\operatorname{Re}(s)<-1,
  \]
  则其拉普拉斯逆变换 $f(t)=$\fillblank。

\item 已知两个离散序列（箭头标明 $n=0$ 的位置）：
  \[
  x[n]=\{3,\,-2,\,\underset{\uparrow}{1},\,0,\,2\},\qquad
  h[n]=\{1,\,\underset{\uparrow}{0},\,-1,\,2\},
  \]
  则卷积和 $y[n]=x[n]*h[n]=$\fillblank。

\item 某无失真传输系统 $H(\mathrm{j}\omega)=2\mathrm{e}^{-\mathrm{j}3\omega}$，激励 $e(t)=3\sin(2t)$，响应 $r(t)=$\fillblank。

\item 如图所示的RLC串联电路，$R=1\,\Omega$，$L=1\,\mathrm{H}$，$C=\dfrac13\,\mathrm{F}$，输入为 $v_{\mathrm{i}}(t)$，输出为电容电压 $v_{\mathrm{C}}(t)$，则系统函数 $H(s)=\dfrac{V_{\mathrm{C}}(s)}{V_{\mathrm{i}}(s)}=$\fillblank。

\begin{center}
\begin{circuitikz}[american,scale=0.9]
  \draw (0,0) to[V, invert, l=$v_{\mathrm{i}}(t)$] (0,3)
              to[R, l=$1\,\Omega$] (3,3)
              to[L, l=$1\,\mathrm{H}$] (5,3)
              to[C, l=$\frac13\,\mathrm{F}$, v=$v_{\mathrm{C}}(t)$] (5,0)
              -- (0,0);
\end{circuitikz}
\end{center}

\item 已知 $X(z)=\dfrac{z}{2(z-0.5)^2}$，收敛域 $|z|>0.5$。设 $y[n]=2^{\,n}\,x[n]$，则 $Y(z)=$\fillblank，收敛域为\fillblank。

\end{enumerate}

% ==================== 三、解答题 ====================
\section*{三、解答题}

\subsection*{1.}

已知LTI连续时间系统的微分方程为
\[
r''(t)+3r'(t)+2r(t)=e'(t)+e(t),
\]
激励 $e(t)=\mathrm{e}^{-3t}u(t)$，系统的全响应为
\[
r(t)=\bigl(\mathrm{e}^{-t}-\mathrm{e}^{-3t}\bigr)\,u(t).
\]

\begin{enumerate}[label=(\arabic*),leftmargin=*]
  \item 指出全响应中的自由响应分量和强迫响应分量；
  \item 利用冲激函数匹配法，求 $r(0_-)$ 和 $r'(0_-)$；
  \item 求系统的单位冲激响应 $h(t)$；
  \item 求系统的零输入响应 $r_{\mathrm{zi}}(t)$ 和零状态响应 $r_{\mathrm{zs}}(t)$。
\end{enumerate}

\subsection*{2.}

某LTI连续时间系统的零极点分布为：极点 $p_1=-1$，$p_2=-4$；零点 $z_1=-2$。已知单位冲激响应的初值 $h(0_+)=1$。

\begin{enumerate}[label=(\arabic*),leftmargin=*]
  \item 求系统函数 $H(s)$，并判断系统是否稳定，说明理由；
  \item 求系统的单位冲激响应 $h(t)$；
  \item 画出系统的直接形式信号流图，并写出输入--输出微分方程；
  \item 若激励 $e(t)=\mathrm{e}^{-2t}u(t)$，求零状态响应 $r_{\mathrm{zs}}(t)$。
\end{enumerate}

\subsection*{3.}

如图所示的频分复用通信系统中，两路基带信号的频谱分别为
\[
E_1(\mathrm{j}\omega)=
\begin{cases}
1, & |\omega|<\omega_0,\\
0, & \text{其他},
\end{cases}
\qquad
E_2(\mathrm{j}\omega)=
\begin{cases}
\sqrt{1-\bigl(\dfrac{\omega}{\omega_0}\bigr)^2}, & |\omega|<\omega_0,\\[8pt]
0, & \text{其他}.
\end{cases}
\]

\begin{center}
\begin{adjustbox}{max width=\textwidth}
\begin{tikzpicture}[
  node distance=0.5cm and 1.0cm,
  block/.style={rectangle,draw,minimum width=1.4cm,minimum height=0.7cm,font=\footnotesize,align=center},
  sum/.style={circle,draw,inner sep=0pt,minimum size=0.55cm,font=\footnotesize},
  dot/.style={circle,fill,inner sep=1pt},
  lbl/.style={font=\footnotesize\color{gray}},
]

% === TRANSMITTER ===
\node (e1) {$e_1(t)$};
\node[sum,right=of e1] (mod1) {$\otimes$};
\node[lbl,above=0.25cm of mod1] {$\cos\omega_1 t$};
\node[sum,right=of mod1] (add) {$+$};
\node[below=0.7cm of add] (e2) {$e_2(t)$};
\node[right=of add] (e3) {$e_3(t)$};
\node[sum,right=of e3] (mod2) {$\otimes$};
\node[lbl,above=0.25cm of mod2] {$\cos\omega_2 t$};
\node[right=of mod2] (e4) {$e_4(t)$};

\draw[->] (e1) -- (mod1);
\draw[->] (mod1) -- (add);
\draw[->] (e2) -- (add);
\draw[->] (add) -- (e3);
\draw[->] (e3) -- (mod2);
\draw[->] (mod2) -- (e4);

% === RECEIVER (below transmitter) ===
\node[sum,below=1.8cm of e4] (demod2) {$\otimes$};
\node[lbl,above=0.25cm of demod2] {$\cos\omega_2 t$};
\node[block,right=of demod2] (h1) {$H_1$\\[-2pt]\tiny 增益$2$, 截止$\omega_0\!+\!\omega_1$};
\node[right=of h1] (e5) {$e_5(t)$};
\node[dot,right=of e5] (split) {};

\draw[->] (e4) -- (e4 |- demod2) -- (demod2);
\draw[->] (demod2) -- (h1);
\draw[->] (h1) -- (e5);
\draw[->] (e5) -- (split);

% Upper branch
\node[block,above right=0.3cm and 0.8cm of split] (h3) {$H_3$\\[-2pt]\tiny 增益$1$};
\node[sum,right=of h3] (mod3) {$\otimes$};
\node[lbl,above=0.25cm of mod3] {$\cos\omega_1 t$};
\node[block,right=of mod3] (h4) {$H_4$};
\node[right=of h4] (out1) {$\frac12 e_1(t)$};

\draw[->] (split) |- (h3);
\draw[->] (h3) -- (mod3);
\draw[->] (mod3) -- (h4);
\draw[->] (h4) -- (out1);

% Lower branch
\node[block,below right=0.3cm and 0.8cm of split] (h2) {$H_2$\\[-2pt]\tiny 增益$1$};
\node[right=of h2] (out2) {$e_2(t)$};

\draw[->] (split) |- (h2);
\draw[->] (h2) -- (out2);

\end{tikzpicture}
\end{adjustbox}
\end{center}

\vspace{0.2cm}

发射端先用 $\cos(\omega_1 t)$ 调制 $e_1(t)$，再与 $e_2(t)$ 叠加得 $e_3(t)$，然后用 $\cos(\omega_2 t)$ 调制得 $e_4(t)$。接收端先用 $\cos(\omega_2 t)$ 解调，经理想低通滤波器 $H_1$（增益为2，截止频率为 $\omega_0+\omega_1$）得 $e_5(t)$。$e_5(t)$ 分两路：上路经 $H_3$（增益为1）、$\cos(\omega_1 t)$ 解调、$H_4$ 后得 $\frac12 e_1(t)$；下路经 $H_2$ 后得 $e_2(t)$。

已知 $\omega_1>2\omega_0$，$\omega_2\gg\omega_1$。

\begin{enumerate}[label=(\arabic*),leftmargin=*]
  \item 画出 $e_3(t)$、$e_4(t)$、$e_5(t)$ 的频谱示意图，标注关键频率值；
  \item 设计 $H_2(\mathrm{j}\omega)$、$H_3(\mathrm{j}\omega)$、$H_4(\mathrm{j}\omega)$ 的频率特性，画出特性曲线并标注关键参数；
  \item 考虑叠加信号 $f(t)=\frac{1}{2}e_1(t)+e_2(t)$。画出其频谱图。若用间隔为 $T_{\!s}$ 的单位冲激序列对 $f(t)$ 进行理想抽样，为保证不失真（不产生频谱混叠），$T_{\!s}$ 的最大值是多少？画出抽样后信号 $f_{\!s}(t)$ 的频谱图。
\end{enumerate}

\subsection*{4.}

某LTI离散时间系统，已知激励 $x[n]=u[n]$ 时的零状态响应为
\[
y[n]=\bigl(8-2\cdot0.5^{\,n}-0.25^{\,n}\bigr)\,u[n].
\]

\begin{enumerate}[label=(\arabic*),leftmargin=*]
  \item 求系统函数 $H(z)$ 及零、极点，判断系统是否稳定；
  \item 求单位样值响应 $h[n]$；
  \item 画出 $|H(\mathrm{e}^{\mathrm{j}\Omega})|$ 的粗略图形，分析系统的频率特性（低通/高通/带通/带阻）；
  \item 若输入 $x[n]=3\cos(n\pi)$，求系统的稳态输出 $y_{\mathrm{ss}}[n]$。
\end{enumerate}

\end{document}
